\subsection{Колебание в трёхмерном пространстве}

\stepcounter{equation}
\begin{equation}\tag{1}\label{eq:wave3-free}
  \begin{cases}
    u_{tt} - a^2 \Delta u = 0, & x \in \R^3,\ t > 0, \\[2pt]
    u|_{t=0} = \varphi(x), \quad u_t|_{t=0} = \psi(x) .
  \end{cases}
\end{equation}

\subsubsection*{Сферические средние}

\begin{definition}
\textbf{Сферическим средним} функции $u$ называется
\[
  \bar u(x,t,z) := \frac{1}{4\pi z^2} \int\limits_{S_{x,z}} u(y,t)\, ds
  = \frac{1}{4\pi z^2} \int\limits_0^{2\pi} \int\limits_0^{\pi}
    u\, z^2 \sin\theta \, d\theta\, d\varphi ,
\]
где $S_{x,z}$ — сфера радиуса $z$ с центром в $x$.
\end{definition}

\begin{warn}
На доске радиальная переменная обозначена буквой, похожей на $z$; это именно
радиус, а не декартова координата. Вторая радиальная переменная, по которой идёт
интегрирование внутри шара, обозначена $\rho$. Здесь сохранены обозначения доски:
$z$ — радиус внешней сферы, $\rho \in [0,z]$ — переменная интегрирования.
\end{warn}

\inputfig{figures/fig-spherical-mean}

\begin{idea}
Множитель $4\pi z^2$ — это площадь сферы $S_{x,z}$, так что $\bar u$ —
действительно среднее по сфере, а не просто интеграл. Отсюда сразу
\[
  \bar u(x,t,z) \xrightarrow[z \to 0]{} u(x,t)
\]
для непрерывной $u$: среднее по стягивающейся сфере стремится к значению в
центре. Именно этот предельный переход в конце даст формулу Кирхгофа.
\end{idea}

\subsubsection*{Уравнение для сферического среднего}

Проинтегрируем уравнение по шару $B_{x,z}$. Левая часть:
\[
  \int\limits_{B_{x,z}} u_{tt}\, dx
  = \int\limits_0^{z} \int\limits_{S_{x,\rho}} u_{tt}\, ds\, d\rho
  = \int\limits_0^{z} 4\pi \rho^2\, \bar u_{tt}\, d\rho .
\]
Правая часть — по первой формуле Грина (она же формула
Гаусса—Остроградского для $\operatorname{div}\nabla u$), с учётом того, что
внешняя нормаль к сфере направлена по радиусу, то есть
$\partial/\partial n = \partial/\partial z$; затем переходим от сферы к
угловым переменным:
\[
  \int\limits_{B_{x,z}} \Delta u\, dx
  = \int\limits_{S_{x,z}} \dd{u}{n}\, ds
  = \int\limits_{S_{x,z}} \dd{u}{z}\, ds
  = \int\limits_0^{2\pi} \int\limits_0^{\pi} \dd{u}{z}\, z^2 \sin\theta\,
    d\theta\, d\varphi
  = z^2 \dd{}{z}\big(4\pi \bar u\big) .
\]
Приравнивая (уравнение даёт $u_{tt} = a^2 \Delta u$) и сокращая на $4\pi$:
\[
  \int\limits_0^{z} \rho^2\, \bar u_{tt}\, d\rho
    - a^2 z^2 \dd{\bar u}{z} = 0 .
\]
Дифференцируем по $z$:
\[
  z^2 \bar u_{tt} - a^2 \dd{}{z}\!\left(z^2 \dd{\bar u}{z}\right) = 0
  \qquad\Longleftrightarrow\qquad
  \bar u_{tt} - a^2 \frac{1}{z^2} \dd{}{z}\!\left(z^2 \dd{\bar u}{z}\right) = 0 .
\]

\begin{idea}
Это уравнение Эйлера—Пуассона—Дарбу: радиальная часть волнового оператора в
$\R^3$. Главное наблюдение следующего шага — в размерности $3$ (и только в
нечётных) оно сводится к \emph{одномерному} волновому уравнению заменой
$v = z\bar u$.
\end{idea}

\subsubsection*{Сведение к полубесконечной струне}

Положим
\[
  \bar u = \frac{v}{z}, \qquad\text{то есть}\qquad v(z,t) = z\,\bar u(x,t,z) .
\]
Тогда
\[
  \begin{cases}
    v_{tt} - a^2 v_{zz} = 0, & z \in (0,\infty), \\[2pt]
    v|_{t=0} = z\,\bar\varphi =: v_0(z), \quad
    v_t|_{t=0} = z\,\bar\psi =: v_1(z), \\[2pt]
    v|_{z=0} = 0 .
  \end{cases}
\]

\begin{idea}
Это ровно задача~\eqref{eq:semi-dirichlet} из \S2: свободные колебания
полубесконечной струны с закреплённым концом. Условие $v|_{z=0} = 0$ не
навязано извне — оно следует из определения $v = z\bar u$ при ограниченной
$\bar u$.

\medskip
Проверка замены: прямое вычисление даёт
$\frac{1}{z^2}\big(z^2 \bar u_z\big)_z = \frac{1}{z}\,(z\bar u)_{zz}$, откуда
после умножения уравнения на $z$ и получается $v_{tt} = a^2 v_{zz}$.
\end{idea}

Здесь $\bar\varphi$ и $\bar\psi$ — сферические средние начальных данных:
\[
  \bar\varphi(x,z) = \frac{1}{4\pi z^2} \int\limits_{S_{x,z}} \varphi(y)\, ds,
  \qquad
  \bar\psi(x,z) = \frac{1}{4\pi z^2} \int\limits_{S_{x,z}} \psi(y)\, ds .
\]

\subsubsection*{Предельный переход и формула Кирхгофа}

Искомое значение восстанавливается как предел среднего:
\[
  u(x,t) = \lim_{z \to 0} \bar u = \lim_{z \to 0} \frac{v(z,t)}{z}
         = \dd{v}{z}\bigg|_{z=0} ,
\]
по правилу Лопиталя: $v|_{z=0} = 0$, так что это неопределённость
$\tfrac{0}{0}$, и предел равен производной $v$ по $z$ в нуле.

Так как $z \to 0$, при фиксированном $t > 0$ мы попадаем в случай $z - at < 0$,
то есть во вторую строку ответа из \S2:
\[
  v = \frac12\Big[v_0(z+at) - v_0(at-z)\Big]
    + \frac{1}{2a}\int\limits_{at-z}^{z+at} v_1(\rho)\, d\rho .
\]
Дифференцируем по $z$ и полагаем $z = 0$:
\[
  \dd{v}{z}\bigg|_{z=0}
  = \frac12\Big[v_0'(at) + v_0'(at)\Big]
    + \frac{1}{2a}\Big[v_1(at) + v_1(at)\Big]
  = v_0'(at) + \frac1a\, v_1(at) .
\]

Остаётся подставить $v_0$ и $v_1$:
\[
  v_0(z) = z\,\bar\varphi = \frac{z}{4\pi z^2} \int\limits_{S_{x,z}} \varphi\, ds
         = \frac{1}{4\pi z} \int\limits_{S_{x,z}} \varphi\, ds ,
\]
\[
  \frac1a\, v_1\big|_{z=at} = \frac1a \big(z\,\bar\psi\big)\Big|_{z=at}
  = \frac1a \cdot \frac{at}{4\pi a^2 t^2} \int\limits_{S_{x,at}} \psi\, ds .
\]
Производная $v_0'$ в точке $z = at$ переписывается через $\partial/\partial t$:
\[
  v_0'(at) = \frac1a\, \dd{}{t}\Big[v_0(at)\Big] ,
\]
так что
\[
  u(x,t) = v_0'(at) + \frac1a v_1(at)
         = \frac1a\left[\dd{}{t} v_0(at) + v_1(at)\right] .
\]

Итого:

\stepcounter{equation}
\begin{equation}\tag{К}\label{eq:kirchhoff-free}
  \boxed{\ u(x,t) = \frac{1}{4\pi a^2}\left[
    \dd{}{t}\left(\frac1t \int\limits_{S_{x,at}} \varphi(y)\, ds_y\right)
    + \frac1t \int\limits_{S_{x,at}} \psi(y)\, ds_y
  \right]\ }
\end{equation}

\begin{center}
\textbf{— формула Кирхгофа} (однородное уравнение).
\end{center}

\begin{idea}
Сравните с компактной записью Даламбера из \S1: структура та же — производная по
$t$ от интеграла по «шару» от $\varphi$ плюс интеграл от $\psi$. Разница
принципиальная: при $n=1$ интегрирование шло по \emph{отрезку}
$\Bxat{x}{at}$, а при $n=3$ — по \emph{сфере} $S_{x,at}$, то есть только по
границе. Отсюда и принцип Гюйгенса (см. ниже).
\end{idea}

\subsubsection*{Неоднородное уравнение: принцип Дюамеля}

Теперь наоборот: однородные начальные условия и неоднородное уравнение. Здесь
\[
  \Box u = u_{tt} - a^2 \Delta u
\]
— \textbf{оператор Даламбера}.

\stepcounter{equation}
\begin{equation}\tag{2}\label{eq:wave3-forced}
  \begin{cases}
    \Box u \equiv u_{tt} - a^2 \Delta u = f(x,t), & x \in \R^3,\ t > 0, \\[2pt]
    u|_{t=0} = 0, \quad u_t|_{t=0} = 0 .
  \end{cases}
\end{equation}

Представим правую часть как суперпозицию мгновенных импульсов:
\[
  f(x,t) = \int\limits_0^{\infty} f(x,\tau)\, \delta(\tau - t)\, d\tau ,
\]
и рассмотрим вспомогательную задачу на один импульс:
\[
  \begin{cases}
    v_{tt} - a^2 \Delta v = f(x,\tau)\,\delta(t - \tau), & t > 0, \\[2pt]
    v|_{t=0} = 0, \quad v_t|_{t=0} = 0 .
  \end{cases}
\]

\paragraph{Что делает импульс.} Проинтегрируем уравнение по
$t \in [\tau-\eps,\, \tau+\eps]$ и устремим $\eps \to 0$:
\[
  \int\limits_{\tau-\eps}^{\tau+\eps} v_{tt}\, dt
  = v_t \Big|_{\tau-\eps}^{\tau+\eps}
  \xrightarrow[\eps \to 0]{} v_t\big|_{t=\tau} ,
  \qquad
  \int\limits_{\tau-\eps}^{\tau+\eps} \Delta v\, dt
  \xrightarrow[\eps \to 0]{} 0 ,
\]
\begin{idea}
Почему в первом пределе остаётся только верхний предел: при $t < \tau$ правая
часть тождественно равна нулю, а начальные данные нулевые, поэтому по теореме
единственности (\S4) там $v \equiv 0$. Значит и $v_t(\tau - \eps) \to 0$.
Во втором пределе всё проще: $\Delta v$ — хорошая функция, а промежуток
интегрирования стягивается в точку.
\end{idea}
\[
  \int\limits_{\tau-\eps}^{\tau+\eps} f(x,\tau)\,\delta(\tau - t)\, dt
  = f(x,\tau) .
\]
Значит $\delta$-импульс не меняет саму $v$, но мгновенно задаёт ей скорость:
\[
  \begin{cases}
    v_{tt} - a^2 \Delta v = 0, & t > \tau, \\[2pt]
    v|_{t=\tau} = 0, \quad v_t|_{t=\tau} = f(x,\tau) .
  \end{cases}
\]

\paragraph{Суперпозиция.} Положим
\[
  u(x,t) = \int\limits_0^{\infty} v(x,t,\tau)\, d\tau
         = \int\limits_0^{t} v\, d\tau
\]
(при $\tau > t$ решение $v$ ещё тождественно равно нулю).

\begin{proposition}
Так построенная $u$ решает задачу~\eqref{eq:wave3-forced}.
\end{proposition}

\begin{proof}
Начальные условия: $u|_{t=0} = 0$ очевидно, и
\[
  u_t = \underbrace{v\big|_{t=\tau}}_{=\,0} + \int\limits_0^t v_t\, d\tau
      = \int\limits_0^t v_t\, d\tau
  \qquad\Longrightarrow\qquad u_t|_{t=0} = 0 .
\]
Далее
\[
  u_{tt} = v_t\big|_{t=\tau} + \int\limits_0^t v_{tt}\, d\tau
         = f(x,t) + \int\limits_0^t v_{tt}\, d\tau ,
\]
и поэтому
\[
  u_{tt} - a^2 \Delta u
  = f(x,t) + \int\limits_0^t \big(v_{tt} - a^2 \Delta v\big)\, d\tau
  = f(x,t) ,
\]
так как подынтегральное выражение равно нулю при $\tau < t$.
\end{proof}

\paragraph{Явный вид.} Задача для $v$
\[
  \begin{cases}
    v_{tt} - a^2 \Delta v = 0, & x \in \R^3,\ t > \tau, \\[2pt]
    v|_{t=\tau} = 0, \quad v_t|_{t=\tau} = f(x,\tau)
  \end{cases}
\]
— это задача~\eqref{eq:wave3-free} со сдвинутым началом отсчёта времени,
$\varphi \equiv 0$ и $\psi = f(\cdot,\tau)$. Применяя
формулу~\eqref{eq:kirchhoff-free} и интегрируя по $\tau$:
\[
  u(x,t) = \frac{1}{4\pi a^2} \int\limits_0^{t} \frac{1}{t-\tau}
  \int\limits_{S_{x,a(t-\tau)}} f(y,\tau)\, ds\, d\tau .
\]

\subsubsection*{Формула Кирхгофа в общем виде}

Складывая оба вклада, для задачи с данными $\varphi, \psi$ и правой частью $f$:
\[
  \boxed{
  \begin{aligned}
    u(x,t) = \frac{1}{4\pi a^2}\Bigg[ &
      \dd{}{t}\left(\frac1t \int\limits_{S_{x,at}} \varphi(y)\, ds_y\right)
      + \frac1t \int\limits_{S_{x,at}} \psi(y)\, ds_y \\[4pt]
    & + \int\limits_0^{t} \frac{1}{t-\tau}
      \int\limits_{S_{x,a(t-\tau)}} f(y,\tau)\, ds_y\, d\tau
    \Bigg]
  \end{aligned}}
\]

\begin{center}
\textbf{— формула Кирхгофа.}
\end{center}

\inputfig{figures/fig-cone-kirchhoff}

\begin{remark}
В тетради к этому рисунку приписано: \emph{область влияния — это только
образующие конуса, в отличие от одномерного случая}. То есть и терминология
здесь та же, что в \S1 («область влияния \emph{на} точку» = область
зависимости), и подчёркнуто главное отличие от $n=1$: работает не весь конус,
а только его боковая поверхность.
\end{remark}

\begin{idea}[title={Принцип Гюйгенса}]
Во все три интеграла входят только \emph{сферы}, то есть боковая поверхность
конуса, а не его внутренность. Поэтому возмущение, сосредоточенное в момент
$t=0$ в малой окрестности точки $y_0$, проходит через точку $x$ за конечное
время $|x-y_0|/a$ и после этого не оставляет следа — сигнал имеет чёткий задний
фронт.

\medskip
При $n=1$ это не так: в формуле Даламбера интеграл от $\psi$ берётся по всему
отрезку $\Bxat{x}{at}$, и «хвост» остаётся навсегда — это видно по ступеньке в
примере с $\psi = \delta$ из \S1.
\end{idea}

\subsubsection*{Замечание: метод спуска и формула Пуассона}

\begin{remark}
Рассмотрим двумерную задачу
\[
  \begin{cases}
    u_{tt} - a^2 \Delta u = f(x,t), & x \in \R^2,\ t > 0, \\[2pt]
    u|_{t=0} = \varphi(x), \quad u_t|_{t=0} = \psi(x) .
  \end{cases}
\]
Она получается из трёхмерной, если считать все данные не зависящими от третьей
координаты:
\[
  f(x_1,x_2,x_3) = f(x_1,x_2) .
\]
Подстановка в формулу Кирхгофа и интегрирование по $x_3$ даёт \textbf{формулу
Пуассона}.
\end{remark}

\begin{warn}
Сама формула Пуассона не выписана ни на доске, ни в тетради. В тетради есть
только приписка: «получим формулу Пуассона — там будет интеграл по шару в
двумерии, проекции трёхмерной сферы». Вывод, по-видимому, в следующей лекции.
\end{warn}

\begin{idea}
Приём называется \textbf{методом спуска}: задача в $\R^2$ решается как задача в
$\R^3$ с данными, постоянными вдоль $x_3$. Интеграл по сфере $S_{x,at} \subset
\R^3$ при этом проецируется на круг $\Bxat{x}{at} \subset \R^2$, причём сфера
накрывает круг дважды — отсюда в формуле Пуассона появляется интеграл по
\emph{кругу} со знаменателем $\sqrt{a^2t^2 - |y-x|^2}$.

\medskip
Принципиальное следствие: в $\R^2$ интегрирование идёт по кругу, а не по
окружности, значит принцип Гюйгенса \textbf{не выполняется} — как и при $n=1$.
\end{idea}
